\chapter{Any Future Metaphysics}
\label{ch:kant}
% ===================================================================

This book opened with a chapter called \emph{Prolegomena to Finding
Mind}.
i borrowed the word from Kant on the first page and then, for three
hundred pages, did not go back for it.
Chapter~\ref{ch:trick} explained why the reclaiming had to wait:
a philosophy that will not show its own method is exactly the thing
the \emph{Prolegomena} was written against, and so the method had
to be shown first.
It has been.
You know now that the whole construction was bought with a
restriction of scope to computation, that the restriction yielded
an ontology and a grounded hypothesis language for free, and that
outside that scope Chapter~\ref{ch:notation} can say what a symbol
system would have to be and not how to get one.

So the borrowed word comes back last, and it comes back in full
view of the trick.
That is a harder position to argue from and a better one to be
believed from.

The Turn is written in the third person that mathematics speaks:
here is a category, here is a fibration, here is a theorem and here
is the hypothesis under which it holds.
Read that way it is a piece of theoretical computer science with
ambitions toward physics, and it can be evaluated as such --- some
of it is proved, some of it is conjectured, and Chapter~\ref{sec:synthesis}
says which is which.
But every construction in it admits a second reading, and the
second reading is what this chapter is for.
The question underneath was never \emph{what is the right
formalism for compressed causal graphs?}
It was: what is it like to be an agent under existential pressure
to work out its own situation?
That is the flip side of the offer artificial intelligence makes
to the sciences.
If useful and salient properties of human intelligence can be
faithfully rendered as computation, then we are entitled to turn
the question around and ask what it is like to be a computation.

Six things the Turn established, which the rest of this chapter
leans on and will not re-derive.
A theory of state and evolution has three pieces --- a grammar, an
equational quotient, a collection of rewrites --- and a theory of
that shape is a graph-structured lambda theory
(Chapter~\ref{ch:gslt}).
When the base rewrites share a distinguished family of head
constructors, the theory is interactive, and the cut between
program and environment becomes perspectival rather than absolute
(Chapter~\ref{ch:igslt}).
Space is not given to such a theory but read off it, as the
structure of cuts through a state, with time the path through the
events those cuts make available (Chapter~\ref{sec:oslf-inputs}).
Alternative carvings of the same dynamics are related by morphisms
that preserve bisimulation rather than syntax, so no carving is
privileged.
Over that base sits the fibration, and an agent's position in it is
a coordinate rather than a rank
(Chapters~\ref{sec:fibration}, \ref{ch:apr}, and~\ref{ch:agy}).
And the agent's hypotheses are formulae of a modal logic generated
from its own substrate and encoded as terms of that substrate,
because no metalanguage is on offer to it
(Chapter~\ref{sec:lts-hml}); testing one costs, the cost is
metered, and the agent does not get it back if it was wrong
(Chapters~\ref{sec:weighted} and~\ref{sec:cost}).

None of that apparatus was postulated.
Each piece is forced by a single existential fact: the full causal
graph of an agent's situation --- every co-occurrence, every
potential influence --- is unworkable as a working representation,
so the agent must compress; and compressed causal graphs have that
shape.
It must compress into rule form, because that is what compression
of state-evolution is.
Its hypotheses must be expressible in its own substrate, because
there is nowhere else to put them.
And they must be testable at a cost commensurate with their reach,
because its resources are finite and its persistence depends on
predictions that pay for themselves.
The structural commitments unfold from the existential one.

Which is a description of a project already attempted.
Kant's transcendental philosophy is a project about what it is like
to be a cognizing subject --- about the conditions of having an
experience at all --- and it starts, as this does, inside the agent
and works outward to what any such agent must be committed to.
Where Kant had to gesture outward, to \emph{us}, his readers, in
our shared cognitive constitution, in order to say what the forms
of cognition are, an agent here can say it from inside its own
resources.
The transcendental project becomes first-personal in a way Kant
could glimpse and could not formalize.

The claim of this chapter, briefly:

\begin{quote}
Kant's transcendental project is vindicated, refined, and
radicalized when its forms of cognition are read as the structural
commitments any computational agent must make in order to predict
its situation, in the only form prediction can take.
\end{quote}

\noindent
The qualifier \emph{computational} in that sentence is not
decoration and is not free, and after Chapter~\ref{ch:trick} nobody
should let me pretend otherwise.
It is the scope restriction, doing its work one last time.
What follows is a reading of Kant for agents whose world is made of
computation.
Whether it is a reading of Kant for us depends on the premise this
book has never stopped stating: that useful and important aspects
of human intelligence can be faithfully rendered as computation.
i think it is the right premise.
i have not shown it, and the difference between thinking and
showing is the whole subject of the two preceding chapters.

\noindent
And the claim underneath that one, which is what the Prestige is
actually for: the whole construction --- a graph-structured lambda
theory, together with the monads and functors this book has
identified on it --- is a proposal for a metaphysics that can be
\emph{built} rather than argued to.
i will come back to that in Section~\ref{sec:constructive-metaphysics},
after the reading of Kant has earned it.

\section{Kant's question}
\label{sec:kant-recap}

The \emph{Prolegomena to Any Future Metaphysics That Will Be Able
to Come Forward as Science} (1783) is Kant's compressed
methodological restatement of the \emph{Critique of Pure
Reason}~\cite{Kant1783,Kant1781}.
The \emph{Critique} had been misread by its reviewers as one more
idealism in the Berkeleyan line --- as the claim that the world is
in some sense mental --- and had thereby failed to communicate its
central architectural innovation.
The \emph{Prolegomena} tries to fix that by being explicit about
method.
Its question is austerely focused: \emph{how is metaphysics
possible as a science?}
And its strategy is regressive.
It begins from the fact that synthetic a priori knowledge exists
--- in mathematics, and in pure natural science --- and reasons
backward to the conditions that must obtain for such knowledge to
be possible at all.

The principal moves, compressed.

\paragraph{The Copernican turn.}
Kant inverts the question of how cognition relates to its objects.
Rather than asking how thought manages to conform to a pre-given
world, he asks how a world of possible experience conforms to
thought.
The forms of intuition --- space and time --- and the categories of
the understanding --- substance, causality, possibility, and the
rest --- are not features of things as they are in themselves.
They are the structure the cognizing subject brings, the conditions
under which anything can be an object of possible experience.

\paragraph{The forms of intuition.}
Space and time are not empirical.
They cannot be extracted from experience, because any experience
already presupposes them as the form in which it is given.
They are pure forms of sensible intuition, a priori in the
strongest sense the word has.

\paragraph{The categories.}
From his table of judgments --- quantity, quality, relation,
modality, each subdivided into three --- Kant derives twelve pure
concepts of the understanding, which structure any judgment about
an object of possible experience.
The principles of pure natural science, that every event has a
cause and that substance persists through change, are the
categories applied to appearances under the schema of time.

\paragraph{The synthetic a priori.}
Synthetic propositions extend knowledge: the predicate is not
contained in the subject.
A priori propositions hold with necessity and universality: their
truth does not wait on experience.
Kant's question is how a proposition can be both.
Mathematics is his exemplar --- $7+5=12$ is, for him, synthetic,
since the concept of twelve is not contained in the concept of
seven-plus-five, and yet a priori, since it holds necessarily.
Pure natural science likewise: the conservation of substance is
synthetic and a priori.
The conditions that make the combination possible are exactly the
forms of intuition and the categories, brought a priori to
experience, and therefore available in advance as knowledge of the
structure of any experience whatever.

\paragraph{Phenomena and noumena.}
The forms and the categories apply to phenomena --- objects as
cognized under those forms.
They do not apply to noumena --- things considered independently of
any cognitive structure.
Noumena are thinkable, in that we cannot rule the concept out.
They are not knowable, because applying the categories beyond
possible experience generates the antinomies: contradictions that
are the mark of reason overreaching.

\paragraph{The limits of metaphysics.}
Speculative metaphysics --- proving the existence of God, the
immortality of the soul, the freedom of the will, from pure reason
alone --- is therefore impossible as a science.
The categories applied outside their domain yield illusion rather
than knowledge.
A future metaphysics, to be a science, must confine itself to the
transcendental: the systematic study of the forms and categories
that constitute the possibility of experience.

That is the architecture, and it is one of the most consequential
achievements of the modern era.
It can be read, point by point, in the structure the Turn built.

\section{The vindication: forms of cognition as constitutive}
\label{sec:kant-vindication}

Kant's deepest move --- the one the framework reproduces with more
precision than he could --- is the claim that the structure of
experience is \emph{constituted} by the cognizing subject's forms
rather than read off things in themselves.
The agent does not meet the world raw and then organize it.
The organization is what makes meeting-the-world possible.

The framework takes the same position with different furniture.
Where Kant has space, time, and twelve categories as fixed
transcendental scaffolding, this book has the GSLT-shape
commitment: the agent's prior that its situation is, at minimum,
some position in the fibration over some interactive
graph-structured lambda theory.
Both are constitutive in the same sense.
The agent literally cannot formulate a hypothesis outside its
scaffolding, because a hypothesis is a term of its own substrate
and the substrate is the scaffolding.
Both are a priori in the relevant sense: not derived from
experience, but presupposed in the having of it.

Then the framework does something Kant explicitly disclaims.
Kant says the forms of intuition and the categories are
\emph{given}.
They are the structure of human cognition; we can describe them,
exhibit them, systematize them, but we cannot derive them from
anything more basic, and he says so.
The framework derives the GSLT-shape commitment from something
more basic, namely that a sensible theory of state and evolution
is a compression of a causal graph.
The shape is not postulated.
It is the form such a compression takes, and the derivation is
short enough to have fit in the opening pages of this chapter.

This is a justification Kant could not have given, and it is worth
being clear about why not, because the reason is not that he was
insufficiently clever.
The resources available to him --- Aristotelian logic, Newtonian
mechanics, Euclidean geometry --- had no way to articulate
compression-of-causal-structure as an underlying notion.
There was no theory of computation to compress into, no notion of a
rewrite rule as a schema standing for a whole family of particular
transitions, no way to say what it would mean for a representation
to be too large to carry.
He had the idea that cognition is constitutive and no vocabulary in
which the constitution could be exhibited as a construction.

So the framework vindicates the transcendental move while replacing
its grounding.
Kant's grounding is empirical and anthropological --- this is how
the human mind works, and we can inspect it --- which is why the
forms come out contingent-looking in a system that needs them
necessary.
The replacement grounding is structural: this is what
theory-formation by a finite computational agent has to look like.
The result is more general than Kant's.
The GSLT-shape commitment binds any computational agent, not only
human ones, and that generality is not an extension of Kant's
result but a consequence of having a better reason for it.

\section{The refinement: the synthetic a priori as compression}
\label{sec:kant-refinement}

Kant's animating problem is the synthetic a priori.
The framework lets it be reread in compression-theoretic terms, and
the rereading is unusually clean.

A synthetic a priori truth is a structural feature of the agent's
prior: a feature of GSLT-shape itself, rather than of any
particular GSLT the agent might go on to hypothesize.
Such a truth is a priori because it follows from the form of the
prior and not from any specific testable claim.
It is synthetic because the prior is substantive --- it constrains
how things can be, it could in principle be wrong --- and not a
tautology.

Take the two analogies Kant leans on hardest.

The Second Analogy, that every event has a cause, becomes the
structural fact that in any GSLT-shaped compression every
state-transition is the firing of some base rule.
There are no uncaused transitions, not because causation has been
proved of the world, but because a transition the agent registers
as a transition is a transition in its theory, and transitions in
its theory are rule firings.
The prior makes this a priori.
Its content --- that every transition the agent can register must
be some rule's firing --- makes it synthetic.

The First Analogy, that substance persists through change, becomes
the structure of Chapter~\ref{sec:oslf-inputs}.
A splitting separates what persists across a rewrite --- the
context $K$, untouched by the firing --- from what changes --- the
subterm $t$ at the cut, consumed and replaced.
The First Analogy is the assertion that the context side of every
cut survives the rewrite at that cut.
And it is worth noticing how thin this is compared with what Kant
wanted from substance.
Nothing here is a bearer of properties, nothing here is the
permanent underlying stuff.
What persists is whatever the rule did not reach.
That is a great deal less than Kant's substance and it is, i think,
exactly as much as anyone was ever entitled to.

The principles of pure mathematics become facts about the freely
generated term type prior to any equational quotient --- the
syntactic backbone the agent must use to formulate any hypothesis
at all, and which is therefore available to it in advance of every
hypothesis.

This is a significant refinement, and the sense in which it is one
is specific.
Kant could never give a satisfying account of \emph{why} the
synthetic a priori principles he identified were the right ones.
The Metaphysical Deduction --- the argument that the table of
judgments yields the table of categories --- is universally
regarded as the weakest link in the \emph{Critique}, and it is
weak in a way that matters, because without it the categories are
a list of things Kant noticed.
The compression reading supplies a non-anthropological reason.
These principles are the structural features of any GSLT-shaped
prior, and any agent holding such a prior must endorse them
whether or not it has ever heard of a table of judgments.

The Turn reached the same conclusion by another road and said so in
passing.
Chapter~\ref{ch:sci} argues that a mortal, interacting learner
cannot afford hypotheses about individuals, because a destructive
assay leaves an individual-level hypothesis worthless after its own
first test, and must therefore traffic in namespaces --- in kinds,
classes, universals.
The traffic in universals that mathematics and science run on is
then neither read off the world nor stipulated by the knower.
It is what such a learner can afford to hold.
That was called, there, a Kantian conclusion with a ledger
attached.
This chapter is what happens when the remark is taken seriously
instead of dropped.

\section{The contestation: the closure of the list}
\label{sec:kant-closure}

Kant insists the table of categories is \emph{complete}.
There are exactly twelve, derivable from four headings of judgment
by a triadic schema, and he treats the completeness as a major
argumentative achievement rather than a happy accident.

The framework cannot endorse the closure.

The structural commitments built into GSLT-shape are far thinner
than twelve categories.
Roughly: grammatical recursion, equational quotient, rewrite as
ground event, interactive structure with distinguished sites.
From these, a good deal of Kant's categorial structure can be
recovered --- substance from grammar, causality from rewrite,
possibility and actuality and necessity from the modal structure of
the logic the theory generates.
But the recovery does not land on twelve, and there is no argument
available that a more articulated list, or a differently
articulated one, would be wrong.

So the framework endorses Kant's project of identifying the
structural commitments built into the cognizing stance while
contesting his closure of it.
The list is open.
The right way to investigate it is not to exhibit a logical schema
and read the categories off its symmetries, but to analyze what
compressed-causal-graph form actually requires --- which is an
ongoing piece of work rather than a finished table, and which this
book has certainly not completed.

i am not sure this counts as a loss.
Kant needed the closure because he needed metaphysics to be
finishable; a science with an open list of first principles was not,
in 1783, a respectable thing to propose.
We have since gotten used to sciences like that.

\section{The reframing: phenomena, noumena, and bisimulation}
\label{sec:kant-noumena}

The framework's most striking parallel to Kant, and the place where
it most clearly exceeds him, concerns the distinction between
phenomena and noumena.

i want to propose the following identification.

\emph{Phenomena are term-level structures.}
They are the agent's actual encounters with situations under its
own apparatus.
Term structure is grammatical, equation-quotiented,
rewrite-animated, differentiable in the sense of
Chapter~\ref{sec:oslf-inputs}; every feature of it is structured by
the agent's GSLT-shape commitment.
Two agents with different presentations of the same dynamics
encounter different term structures --- exactly as two of Kant's
subjects with different forms of intuition would encounter
different phenomena.

\emph{Noumena are bisimulation equivalence classes.}
They are what survives the variation across presentations.
Two terms living in different GSLTs but in corresponding
bisimulation classes are the same in the sense this book takes
seriously, even though no agent at any particular position in the
fibration cognizes that identity directly.
It can only posit it, by recognizing that the morphisms of the
semantic category preserve bisimulation.
The bisimulation class is what the agent reasons toward, not what
it encounters; it is an invariant that transcends any particular
carving.

The identification is faithful to Kant in three respects, and it is
worth being explicit about them, because a parallel that is merely
suggestive is not worth drawing.
Phenomena are \emph{structured} --- already categorized, already
spatiotemporal, already objects of experience --- rather than raw
sense data.
Noumena are \emph{invariant} under variation of the categorial
structure: they are what would remain if one could think the
situation independently of any particular carving.
And the agent cannot know noumena directly; it can only posit them
as the substrate-independent invariants that its morphisms must
preserve.

Where the framework exceeds Kant is in what can then be said about
them.

Kant refuses positive characterization of the noumenal because he
has no apparatus by which a subject's own conceptual resources
could reach past its own categorial carving and remain
intelligible.
He distinguishes the \emph{negative} concept of the noumenon ---
the object considered as not-categorized --- from a \emph{positive}
one --- the object as it is in itself --- and withholds confidence
from the positive concept for the rest of his career.

Hennessy--Milner adequacy is exactly the apparatus he lacked.
Two states are bisimilar if and only if they satisfy the same
formulae of the modal logic.
The bisimulation class is therefore \emph{positively
characterized}, and characterized by the very logic the agent's
hypotheses are already written in.
The noumenal can be reached --- not as raw thing-in-itself, but as
the equivalence class picked out by the modal logic under variation
of presentation.
And because the logic is generated functorially and preserved by
the morphisms, a formula characterizing a class at one point in the
fibration translates to a formula characterizing the corresponding
class at another.
The agent's conceptual apparatus travels with it.
Its hypotheses keep their meaning across substrate changes --- which
they had better, since the agent cannot rule out that what it took
for its substrate is only a presentation, and that the genuine
invariants of its situation live at the level of bisimulation
rather than at the level of syntax.

Notice what this does to the shape of the distinction.
In Kant, phenomenon and noumenon are \emph{vertical}: two aspects of
the same object, the cognized aspect and the in-itself aspect,
stacked.
Here the analog is \emph{horizontal}, running across the semantic
category: term structures in different theories are different
phenomenal presentations, and the bisimulation class is the
equivalence class running across them.
That is closer to a gauge-theoretic picture than to a layered one.
The noumenal is the gauge-invariant content; the phenomenal is the
gauge-dependent presentation; different agents at different points
in the fibration are working in different gauges; and the
bisimulation class is what they all agree on.

It also dissolves what is usually taken to be Kant's hardest
problem, the apparent tension between \emph{we cannot know things
in themselves} and \emph{we are committed to their existence}.
The agent is not committed to a layer of reality it cannot access.
It is committed to the form of its hypotheses, and that form
positively determines the substrate-independent invariants.
The antinomies come out reframed as well.
They are what happens when the agent applies its GSLT-shape
commitment outside the regime where compression into causal-graph
form is reliable.
They are diagnostic of prior-misapplication rather than of
metaphysical paradox --- which is a considerably less romantic
diagnosis, and a more actionable one.

\subsection{Characterized, and unaffordable}
\label{sec:kant-affordable}

Two chapters of the Turn appear to say the opposite of all this,
and the appearance is worth taking seriously rather than
smoothing over, because the resolution is the pivot of the entire
book.

Chapter~\ref{sec:false-ontology} argues that an agent embedded in a
coherence cluster will meet the cluster hierarchy first, will build
excellent theories of it, and will exhaust its budget long before
the bisimulation quotient becomes experimentally accessible.
It will mistake the hierarchy for the whole story, and this is not
a failure of its method but a consequence of the hierarchy being
the dominant accessible phenomenon at its scale.
Chapter~\ref{ch:sci} puts a ledger under the same point:
Proposition~\ref{sci:prop:limit} says that the complete theory of an
environment sits at the top of the lattice as a limit the scientist
approaches and never reaches, and that no learner with a finite
token supply can install it.

So which is it?
Is the noumenal reachable or is it not?

Both, and the distinction is the one worth having.

Hennessy--Milner adequacy is a statement about \emph{characterization}.
The bisimulation class is picked out, with no remainder, by
formulae of the agent's own logic.
Nothing in the theorem says the agent can afford to evaluate them.
The formulae that separate two behaviors may be arbitrarily deep,
each additional modal depth costs, and the cost is metered against
a budget the agent needs for other things --- staying alive, mostly.
The agent's logic determines the noumenal without the agent being
able to buy it.

This is Kant's \emph{thinkable but not knowable}, and it arrives
translated.
For Kant the barrier is categorial: the categories do not apply
beyond possible experience, so the attempt to know the noumenal is
not expensive but malformed, and produces contradiction rather than
expense.
Here the barrier is a budget.
The attempt is perfectly well-formed.
It is simply not affordable, by this agent, at this position, with
this much left in the account.

And a budget constraint is a different kind of thing from a
categorial prohibition.
It can be relaxed.
It can be financed --- Chapter~\ref{ch:gam} is in large part about
what it costs to buy another layer of resolution and when the
purchase pays for itself.
It can be pooled: Chapter~\ref{ch:com} composes learners into
populations that can collectively afford distinctions no member can
afford alone, which is a tolerable formal description of what a
scientific community is for.
It can be moved: an agent at a different position in the fibration
has different things cheaply available to it, and the tower is,
among other things, a ranking by what is affordable.
None of these dissolve the limit.
Proposition~\ref{sci:prop:limit} is not repealed by any amount of
money.
But they change the character of the situation from a wall to a
frontier, and the difference between a wall and a frontier is the
difference between a metaphysics that ends in resignation and one
that ends in a research program.

There is a third thing that can be done to a budget constraint, and
it is the one that gives the abstraction a number.
It can be \emph{measured}.
Chapter~\ref{sec:scope-manufacture} carries out the same argument
one level down, about a question so mundane that nobody would have
expected metaphysics to be involved: how many of something there
are.
A count needs a region to count in, and the region a learner has is
the part of the world it can pay to visit; a count needs a criterion
for sameness, and the criterion a learner has is the depth of
distinction it was assembled to afford.
Both are budget constraints, and both are computable.
What comes out is that a learner does not fail to see the
distinctions it cannot afford --- it \emph{fails to see that they
are there}, and reports a world with fewer kinds in it than the
world has.

So the noumenal has a mundane cousin, and the cousin can be
quantified. What is out of reach is not a separate realm; it is the
part of this one that lies past the last distinction the budget
would buy, and how far past can be worked out from what the learner
is made of.

There is a third instance, and Chapter~\ref{ch:notation} has
already supplied it, though it was doing something else at the time
and did not press the point.
There is now a fourth as well, which arrived after this section was
written and which i have left where it lands rather than folding in:
Chapter~\ref{ch:sim}'s descending case, where what a richer world
distinguishes is unavailable here not by decree but because no
inhabitant can pay the choice strength that resolving it would take.
Remark~\ref{rmk:sim-noumenon} makes the point; it belongs in the
gathering below and the reader should count it there.
The first was about which \emph{questions} an agent can afford:
modal depth, and the separating formula it cannot buy.
The second was about how many \emph{things} it can afford to
distinguish: a count, and the kinds it does not know are missing.
The third is about which marks count as the same mark --- whether
the agent has a symbol system at all in any exact sense, or only a
smear.
Goodman's requirement of finite differentiation is a separation
condition that behavioral equivalence, being a greatest fixed
point, is systematically silent about;
Conjecture~\ref{conj:fd} says a budget supplies it, with the budget
itself as the separation constant.
If that is right then notational adequacy is purchasable too, a
poorer agent has a coarser notation in a precise sense, and the
line between the digital and the analogue --- treated as a fact
about kinds of inscription since Goodman drew it --- turns out to
be a fact about what an observer can pay for.

Three instances is enough to stop calling it a coincidence.
What is going on is that every apparent limit on what a
computational agent can know, once you look at it closely enough to
say where the limit is, resolves into a price.
Not one of them resolves into a prohibition.
That is the pivot, and it is the single most Kantian and least
Kantian thing in this book at the same time: Kant's conclusion
survives the translation into computation --- the thing in itself
really is out of reach --- but it survives as an economic fact
rather than a logical one, and economic facts have the property
that you can do something about them, at a price, up to a point
that can be calculated.

\section{The radicalization: predictive existence}
\label{sec:kant-radical}

Kant's transcendental subject is, structurally, static.
The forms of intuition and the categories are fixed.
They do not evolve; they do not compete with alternatives; there is
no sense in which the subject's commitment to them is at stake.
The transcendental subject simply \emph{has} these forms, and the
question of whether having them was a good idea does not arise and
could not be posed.

The framework introduces a dimension Kant has no room for.

The agent's commitment to GSLT-shape is existentially loaded.
Testing a hypothesis costs --- the source program for this book
called the currency phlogiston, phlo, for reasons that were
accidental and have stuck --- and the metering is not decorative.
Correct prediction recovers most of what was spent.
Incorrect prediction does not.
The balance accumulates, and it determines whether the agent goes
on computing at all.
The forms are not merely had.
They are paid for, and the payments are calibrated by predictive
success.

Three consequences, each of which Kant would have found
objectionable and each of which follows.

\paragraph{Transcendental philosophy becomes empirical.}
The question \emph{is the GSLT-shape commitment correct?} is not
answered a priori.
It is answered by whether agents holding it persist.
This is a pragmatist turn, or an evolutionary one, and Kant
explicitly rejects both; but the framework forces it, because the
agent's epistemic situation simply is its existential situation and
there is no third place to stand from which to separate them.

\paragraph{The unity of apperception acquires a thermodynamic reading.}
Kant's \emph{I think} must be able to accompany all of the
subject's representations; the synthetic unity of consciousness is,
for him, a structural condition, something the subject has by being
a subject.
Here, the agent's continued capacity to formulate and test
hypotheses is what constitutes its persistence as an agent.
The unity is a dynamic equilibrium: it holds as long as
hypothesis-formation pays for itself.
Lose enough, and the capacity for synthetic unity dissolves --- not
as a metaphor, and not as a philosophical embarrassment, but as a
computational fact with a second law behind it
(Chapter~\ref{sec:cost}).

\paragraph{The transcendental subject becomes a population.}
Different agents can hold different priors: different positions in
the fibration, different families of interaction sites, different
semirings on the decorations.
All are under the same selection pressure.
``The transcendental subject'' is therefore not a singular entity
but a distribution over possible priors, shaped by which priors
survive in which environments --- which is exactly the object
Chapter~\ref{ch:com} constructs, and exactly the reason it
constructs it.
That there is a non-trivial distribution of viable priors at all is
a substantive fact about reality.
A world in which only one prior worked would be a world with a
transcendental subject in Kant's singular sense; a world in which
none did would have no subjects.

One caution about the third of these.
It is tempting, having noticed that decorations can take values in
any semiring, to say that real-valued weights give classical
physics and complex-valued weights give quantum physics, and to add
quantum agents to the population as a further variety.
Chapter~\ref{sec:weighted} does establish the semiring-parametric
family, and complex weights are perfectly admissible decorations.
But Chapter~\ref{sec:quantum} is where that hope goes to be
examined, and it comes back with an obstruction rather than a
result: bisimulation refuses to pool the runs that would have to be
pooled, and the only normalization under which the complex reading
lifts the real one has already added them where nothing can cancel.
What it also comes back with is the location of the one slot a
presentation does have for coherence --- the equations it declines
to quotient --- which is a real finding and is not the one that was
hoped for.
What varies across the population is the semiring.
Whether varying it that far buys quantum mechanics is open, and the
chapter that says so also says why it is unfinished.

\section{The methodological inversion}
\label{sec:kant-method}

Kant's method is regressive.
He takes the existence of synthetic a priori knowledge as given and
reasons backward to its conditions, and he defends the choice
explicitly as appropriate for readers who already grant the fact
and want the explanation.

The framework permits a constructive method that Kant thought
impossible.
Rather than presupposing the success of synthetic a priori
knowledge and reasoning to conditions, it begins from the structure
of GSLT-shape and derives what an agent committed to that shape
must know.
The synthetic a priori is not a starting datum.
It is a consequence of the form of the prior.

This unifies what Kant had to keep apart.
The framework gives a single account that runs regressively --- the
conditions of the agent's persistence are the shape commitment and
its position in the fibration --- and progressively --- from the
shape commitment, derive the structural commitments the agent must
hold.
The two directions meet, and they meet at the same place, which is
the sort of thing that suggests the place is real.

\section{A constructive metaphysics}
\label{sec:constructive-metaphysics}

Which brings me to the claim i deferred at the beginning.

Kant called his enterprise a \emph{propaedeutic} --- a
preparation, a clearing of ground, a specification of what a future
metaphysics would have to look like if it were ever to arrive.
He was careful about the word.
The \emph{Prolegomena} does not offer a metaphysics.
It offers the conditions any metaphysics must satisfy, and then
stops, because the resources for going further did not exist.

This book's proposal is that the going-further can now be
attempted, and that what it consists of is a construction rather
than an argument.

The construction is this.
A graph-structured lambda theory: a grammar, an equational
quotient, a set of rewrites.
On it, the reversibility construction of
Chapter~\ref{sec:reversibility}, which attaches a recording device
to a state and gets the reversible time of physics.
The history endofunctor and the cost monad of
Chapters~\ref{sec:weighted} and~\ref{sec:cost}, which meter what
happens and pay for it, and whose interaction is the difference
between a ledger that balances and one that
does not (Remark~\ref{rmk:twomonads}).
The derivative, which manufactures space out of term structure,
and the event sub-bundle, which restricts it to what can actually
happen (Chapter~\ref{sec:oslf-inputs}).
The OSLF functor, which manufactures the agent's logic out of the
theory's own operational semantics rather than stipulating it
alongside.
The fibration, which lets the agent's own computational strength be
a coordinate it occupies rather than a fact it must assume.
And the semantic category, whose morphisms preserve bisimulation
and thereby say which of all this was presentation and which was
content.

Each of these is a functor, a monad, or a fibration --- something
with a construction, a universal property, and the obligation to
compose with its neighbors.
None of them is a posit.
Together they are what i want to offer as a \emph{constructive
metaphysics}: not a set of claims about what there fundamentally
is, defended by argument against rival sets of claims, but an
apparatus that, given a theory of state and evolution, builds out
of it a space, a time, a history, a cost, a logic, an agent, and a
notion of what that agent cannot know and what it would cost to
find out.

The difference between this and a metaphysics in the traditional
sense is the difference between a proof and an assertion.
The traditional question --- \emph{what is there?} --- gets
answered here in a shape that will look evasive at first and is
not: there is whatever a compressed causal graph has in it, and the
apparatus tells you what that is, given the graph.
Substance is what the rewrite did not reach.
Causation is a rule firing.
Space is a structure of cuts.
Time is a path.
The knowable is what the budget covers; the real is what survives
the change of presentation.
Each of those is a construction with a definition behind it, and
each of them can be wrong in a way that would show.

And the parts that are not yet built are visible as gaps in a
construction rather than as disagreements between schools.
The OSLF functor is described in this book and not exhibited.
The metric that would make behavioral nearness into actual nearness
is not yet derived.
The complex case is obstructed.
These are not positions to be defended.
They are things to be finished, and the difference is most of what
i mean by calling the proposal constructive.

\section{Coda: transcendental philosophy made internal}
\label{sec:kant-coda}

The deepest thing the framework does to Kant, on this reading, is to
make transcendental philosophy \emph{internal} to the kind of agent
whose cognition it describes.

Kant had to gesture outward.
To articulate the forms of our cognition he had to address us --- his
readers, in our shared human constitution --- because there was no
way for a subject to get at its own forms from inside.
The forms are the condition of experience and therefore not among
the things experienced, and so the transcendental philosopher has
to stand outside the subject to describe what the subject is
standing in.
That gesture is doing a lot of work in Kant, and it is fragile: it
is why the whole enterprise reads as anthropology to its critics.

The framework needs no such gesture.
An agent inside it can, with sufficient resources at its position,
formulate and test hypotheses about its own shape commitment,
because its hypotheses are terms of the same substrate the
commitment is about.
The reflection is not a philosophical posture but a feature of the
calculus --- which is, after all, what a reflective higher-order
calculus was built to have.
Transcendental philosophy becomes a first-person, existentially
loaded, computational practice.

What Kant called the propaedeutic to all future metaphysics, the
framework can call the structure of any agent that endures by
predicting.

Space emerges from term structure.
Time emerges from paths through it.
History emerges from a recording device carried along.
Phenomena are what the agent encounters under its own carving;
noumena are the substrate-independent invariants its logic can
characterize and its budget cannot always buy.
The transcendental subject is a viable position in the fibration,
and there are many of them.
The synthetic a priori is the structural content of a compressed
causal graph.

All of it follows from one commitment, which is the commitment the
Pledge asked you to take seriously three hundred pages ago: that
useful and important aspects of human intelligence can be
faithfully rendered as computation.
If they can, then cognition requires at least as much structure as
models of computation do.
Models of computation are compressed causal graphs.
Compressed causal graphs are GSLTs.
And from GSLTs, equipped as above, one can generate notions of
space and time rich enough to recover a great deal of physics.
Doing it in that order --- asking first how much structure a
reasoning agent needs, and only then deriving what can be derived
of physics from that structure --- solves several problems at once,
and one of them is the problem of why the physics should have been
knowable by anything in the first place.

Kant would, i think, have recognized the architecture.
He could not have anticipated the resources by which it could be
filled in.
But the \emph{Prolegomena} asked the right question, and by a
sufficiently long way around --- through Church and Turing, through
Hennessy and Milner and Plotkin, through Huet and McBride and the
process calculus tradition --- we are in a position to start
answering it.

One more thing has to be said before that, though, and it is not a
construction.
Everything in this chapter --- the vindication, the refinement, the
reframing of phenomena and noumena, the whole first-personal
transcendental project --- rests on the agent's hypotheses being
terms of the same substrate its hypotheses are about.
That is the condition under which Kant's outward gesture becomes
unnecessary, and it is the condition that made all of this
available.
It is also a condition we arranged, and
Chapter~\ref{ch:trick} said how.
i have put the reading of Kant last rather than first so that it
would have to stand up in that light, and i think it does.
What it is a reading of, precisely, is what it is like to be an
agent whose world is made of computation.
Whether that is a reading of us is the premise, and the premise is
still a premise.

After that:
the next thing that has to be built is a metric --- a notion of
distance under which computational behaviors that are nearly the
same come out near each other.
With that in hand, the space read off a term structure stops being
merely combinatorial and starts being the kind of space one can do
physics in, and the derivation runs the whole way --- from causal
models rich enough to support cognition, to world models, rather
than the other way around.
That is the next problem, and this book has not solved it.
